\section{General Syntax and Semantics}
    % \epigraph{If I had the chance, I would go back to 10 September, 2001 to warn people about Covid-19.}{\textit{Donald J. Trump}}
\subsection{Primitive Data Types}
\subsubsection{Integers}
AFT provides a complete set of integer types with explicit bit-width specifications. \\\\
\textbf{Unsigned Integers:}
\begin{table}[H]
    \centering
    \begin{tabular}{cccc}
        \hline
        \textbf{Type} & \textbf{Size} & \textbf{Minimum Value} & \textbf{Maximum Value} \\
        \hline
        \mintinline{rust}{u8} & 8-bit & 0 & 255 \\
        \mintinline{rust}{u16} & 16-bit & 0 & 65535 \\
        \mintinline{rust}{u32} & 32-bit & 0 & $2^{32} - 1$ \\
        \mintinline{rust}{u64} & 64-bit & 0 & $2^{64} - 1$ \\
        \mintinline{rust}{u128} & 128-bit & 0 & $2^{128} - 1$ \\
        \hline
    \end{tabular}
\end{table}
\noindent{\textbf{Signed Integers:}}
\begin{table}[H]
    \centering
    \begin{tabular}{cccc}
        \hline
        \textbf{Type} & \textbf{Size} & \textbf{Minimum Value} & \textbf{Maximum Value} \\
        \hline
        \mintinline{rust}{i8} & 8-bit & -128 & 127 \\
        \mintinline{rust}{i16} & 16-bit & -32768 & 32767 \\
        \mintinline{rust}{i32} & 32-bit & $-2^{31}$ & $2^{31} - 1$ \\
        \mintinline{rust}{i64} & 64-bit & $-2^{63}$ & $2^{63} - 1$ \\
        \mintinline{rust}{i128} & 128-bit & $-2^{127}$ & $2^{127} - 1$ \\
        \hline
    \end{tabular}
\end{table}
By default, any integer literals without an explicitly mentioned type will be inferred as \mintinline{rust}{i64}.

\subsubsection{Floating-Point Numbers}
AFT provides IEEE 754-compliant floating-point types for real number computations, with support for both single and double precision arithmetic.
\begin{table}[H]
    \centering
    \begin{tabular}{ccccc}
        \hline
        \textbf{Type} & \textbf{Size} & \textbf{Mantissa} & \textbf{Min. Exp.} & \textbf{Max. Exp.} \\
        \hline
        \mintinline{rust}{f32} & 32-bit & 24-bit & -126 & 127 \\
        \mintinline{rust}{f64} & 64-bit & 53-bit & -1022 & 1023 \\
        \hline
    \end{tabular}
\end{table}
By default, any floating-point literals without an explicitly mentioned type will be inferred as \mintinline{rust}{f64}.

\subsubsection{Complex Numbers}
AFT provides native support for complex numbers, essential for signal processing. Complex numbers are represented as pairs of floating-point values for the real and imaginary components.

\begin{table}[H]
    \centering
    \begin{tabular}{ccc}
        \toprule
        \textbf{Type} & \textbf{Real Part} & \textbf{Imaginary Part} \\
        \midrule
        \mintinline{rust}{c32} & \mintinline{rust}{f32} & \mintinline{rust}{f32} \\
        \mintinline{rust}{c64} & \mintinline{rust}{f64} & \mintinline{rust}{f64} \\
        \bottomrule
    \end{tabular}
\end{table}

By default, complex number literals without explicit typing are inferred as \mintinline{rust}{c64}.

\subsubsection{Boolean}
AFT provides a standard boolean type for logical operations and control flow.

\textbf{Type:} \mintinline{rust}{bool}

\subsubsection{Strings}
AFT provides a standard type for strings. Strings accept all valid ASCII characters. No provisions have been made to work with Unicode characters or characters from other character sets.
\begin{table}[]
    \centering
    \begin{tabular}{c|c}
         &  \\
         & 
    \end{tabular}
    \caption{Caption}
    \label{tab:placeholder}
\end{table}

\subsubsection{Vectors}
AFT provides built-in support for homogeneous collections through vector types, enabling efficient storage and manipulation of sequences of elements.

\textbf{Type:} \mintinline{rust}{[T]} \\
Vectors can also be initialized using ranges:
\begin{minted}{rust}
let x = a..b; // [ a, a + 1, a + 2, ..., b - 1 ]
let x = a..=b; // [ a, a + 1, a + 2, ..., b ]
\end{minted}
\subsection{Keywords}
AFT reserves the following keywords for language constructs:

\begin{table}[H]
    \centering
    \begin{tabular}{cccc}
        \toprule
        \mintinline{rust}{let} & \mintinline{rust}{const} & \mintinline{rust}{fn} & \mintinline{rust}{return} \\
        \mintinline{rust}{if} & \mintinline{rust}{else} & \mintinline{rust}{while} & \mintinline{rust}{for} \\
        \mintinline{rust}{repeat} & \mintinline{rust}{struct} & \mintinline{rust}{enum} & \mintinline{rust}{true} \\
        \mintinline{rust}{false} & \mintinline{rust}{as} & \mintinline{rust}{i8} & \mintinline{rust}{u8} \\
        \mintinline{rust}{i16} & \mintinline{rust}{u16} & \mintinline{rust}{i32} & \mintinline{rust}{u32} \\
        \mintinline{rust}{i64} & \mintinline{rust}{u64} & \mintinline{rust}{i128} & \mintinline{rust}{u128} \\
        \mintinline{rust}{f32} & \mintinline{rust}{f64} & \mintinline{rust}{c32} & \mintinline{rust}{c64} \\
        \mintinline{rust}{bool} & \mintinline{rust}{str} & \mintinline{rust}{continue} & \mintinline{rust}{break} \\
        \mintinline{rust}{return} & \mintinline{rust}{static} & \mintinline{text}{PI} \\
        \bottomrule
    \end{tabular}
\end{table}

\subsection{Literals}
\subsubsection{Boolean Literals}
\begin{itemize}
    \item \mintinline{rust}{true}
    \item \mintinline{rust}{false}
\end{itemize}

\subsubsection{Integer Literals}
\begin{itemize}
    \item \textbf{Decimal (default):} \mintinline{rust}{42}, \mintinline{rust}{123}
    \item \textbf{Binary:} \mintinline{rust}{0b1010}, \mintinline{rust}{0b11110000}
    \item \textbf{Octal:} \mintinline{rust}{0o755}, \mintinline{rust}{0o644}
    \item \textbf{Hexadecimal:} \mintinline{rust}{0xFF}, \mintinline{rust}{0xDEADBEEF}
\end{itemize}

\subsubsection{Float Literals}
The integer part can be omitted if it is 0. The decimal part can only be omitted if it is zero and an exponent is present.
\begin{itemize}
    \item \mintinline{rust}{PI} - 3.14149... , pi is available as float const 
    \item \mintinline{rust}{.5} - Integer part omitted
    \item \mintinline{rust}{2.0}
    \item \mintinline{rust}{0.5}
    \item \mintinline{rust}{2.5e10} - Scientific notation
    \item \mintinline{rust}{1e6} - Scientific notation with omitted decimal part
    \item \mintinline{rust}{.3e-6} - Scientific notation with omitted integer part
    \item Various ways to write zero: \mintinline{text}{.0, 0.0}
\end{itemize}

\subsubsection{Complex Literals}
The imaginary part needs to be explicitly specified for the literal to be inferred as complex.
\begin{itemize}
    \item \mintinline{rust}{3.0 + 4.0j} - Real and imaginary parts
    \item \mintinline{rust}{2.5j} - Pure imaginary number
    \item \mintinline{rust}{0+0j} - Complex zero
    \item \mintinline{rust}{6 - 9j} - Negative imaginary part
    \item \mintinline{rust}{6 - j} -  imaginary part without constant
\end{itemize}

\subsubsection{Vector Literals}
Vector literals include comma-separated elements enclosed in square brackets.
\begin{itemize}
    \item \mintinline{rust}{[1, 2, 3]} - Integer vector
    \item \mintinline{rust}{[1.0, 2.0, 3.0]} - Float vector
    \item \mintinline{rust}{[1, 2j, -3, -4j]} - Complex vector
\end{itemize}

\subsection{String Literals}
String literals must be enclosed in double quotes \mintinline{rust}{"}. AFT supports the following escape sequences in string literals:

\begin{table}[h]
    \centering
    \begin{tabular}{cc}
        Escape Sequence & Character Represented \\
        \hline
        \textbackslash n & Line Feed (LF) \\
        \textbackslash r & Carriage Return (CR) \\
        \textbackslash t & Horizontal Tab \\
        \textbackslash v & Vertical Tab \\
        \textbackslash\textbackslash & Bacslash \\
        \textbackslash" & Double Quotation Mark \\
        \hline
    \end{tabular}
\end{table}

\begin{minted}{rust}
// This is a sample "string".
let s = "This is a sample \"string\".";
\end{minted}

\subsection{Identifiers}
Identifiers in AFT are case-sensitive and must follow these rules:
\begin{itemize}
    \item Must start with an alphabetic character or underscore
    \item Can contain alphanumeric characters and underscores
    \item Cannot be a reserved keyword
\end{itemize}

Valid identifiers: \mintinline{rust}{variable_name}, \mintinline{rust}{_private}, \mintinline{rust}{camelCase}, \mintinline{rust}{CONSTANT_VALUE}

\subsection{Operators}
\subsubsection{Unary Operators}
\begin{table}[H]
    \centering
    \begin{tabular}{cc}
        \toprule
        \textbf{Operator} & \textbf{Description} \\
        \midrule
        \mintinline{rust}{+} & positive \\
        \mintinline{rust}{-} & negative \\
        \bottomrule
    \end{tabular}
\end{table}
\subsubsection{Arithmetic Operators}
\begin{table}[H]
    \centering
    \begin{tabular}{cc}
        \toprule
        \textbf{Operator} & \textbf{Description} \\
        \midrule
        \mintinline{rust}{+} & Addition \\
        \mintinline{rust}{-} & Subtraction \\
        \mintinline{rust}{*} & Multiplication \\
        \mintinline{rust}{/} & Division \\
        \mintinline{rust}{%} & Modulo \\
        \mintinline{rust}{**} & Exponentiation \\
        \bottomrule
    \end{tabular}
\end{table}

\subsubsection{Comparison Operators}
\begin{table}[H]
    \centering
    \begin{tabular}{cc}
        \toprule
        \textbf{Operator} & \textbf{Description} \\
        \midrule
        \mintinline{rust}{==} & Equal to \\
        \mintinline{rust}{!=} & Not equal to \\
        \mintinline{rust}{<} & Less than \\
        \mintinline{rust}{>} & Greater than \\
        \mintinline{rust}{<=} & Less than or equal to \\
        \mintinline{rust}{>=} & Greater than or equal to \\
        \bottomrule
    \end{tabular}
\end{table}

\subsubsection{Logical Operators}
All the logical operators are short-circuit operators. They can only be applied to values of type \mintinline{rust}{bool}.
\begin{table}[H]
    \centering
    \begin{tabular}{cc}
        \toprule
        \textbf{Operator} & \textbf{Description} \\
        \midrule
        \mintinline{rust}{&&} & Logical AND \\
        \mintinline{rust}{||} & Logical OR \\
        \mintinline{rust}{^} & Logical XOR \\
        \mintinline{rust}{!} & Logical NOT \\
        \bottomrule
    \end{tabular}
\end{table}

\subsubsection{Bitwise Operators}
AFT offers standard bitwise operators that can be applied to any integer type. Both operands of binary bitwise operators must be of the same type.
\begin{table}[H]
    \centering
    \begin{tabular}{cc}
        \toprule
        \textbf{Operator} & \textbf{Description} \\
        \midrule
        \mintinline{rust}{&} & Bitwise AND \\
        \mintinline{rust}{|} & Bitwise OR \\
        \mintinline{rust}{^} & Bitwise XOR \\
        \mintinline{rust}{<<} & Bitwise left shift \\
        \mintinline{rust}{>>} & Bitwise right shift \\
        \mintinline{rust}{!} & Bitwise NOT \\
        \bottomrule
    \end{tabular}
\end{table}

\subsubsection{String Operators}
\begin{table}[H]
    \centering
    \begin{tabular}{cc}
        \hline
        \textbf{Operator} & \textbf{Description} \\
        \hline
        \mintinline{text}{+} & Concatenation \\
        \mintinline{text}{[]} & Subscripting \\
        \hline
    \end{tabular}
\end{table}
\textbf{Note:} Strings are 0-indexed.

\subsubsection{Vector Operators}
\begin{table}[H]
    \centering
    \begin{tabular}{cc}
        \toprule
        \textbf{Operator} & \textbf{Description} \\
        \midrule
        \mintinline{text}{||} & Concatenation \\
        \mintinline{text}{.*} & Convolution \\
        \mintinline{text}{+} & Vector Addition (or Unary) \\
        \mintinline{text}{-} & Vector Subtraction (or Unary) \\
        \mintinline{text}{*} & Scalar Multiplication \\
        \mintinline{text}{/} & Scalar Division \\
        \mintinline{text}{[]} & Subscripting \\
        \mintinline{text}{~} & Reversing the order of elements \\
        \bottomrule
    \end{tabular}
\end{table}
\textbf{Note:} Vectors are 0-indexed.

\subsubsection{Type Casting with \texttt{as}}

In AFT, explicit type casting is done using the \mintinline{text}{as} keyword .

\begin{table}[H]
    \centering
    \begin{tabular}{cc}
        \toprule
        \textbf{Expression} & \textbf{Description} \\
        \midrule
        \mintinline{rust}{x as i32} & Casts \mintinline{text}{x} to integer \\
        \mintinline{rust}{y as f32} & Casts \mintinline{text}{y} to floating-point \\
        \mintinline{rust}{z as c64} & Casts \mintinline{text}{z} to complex number \\
        \mintinline{rust}{v as [c64]} & Converts vector \mintinline{text}{v} into a vector of \mintinline{rust}{c64} type \\
        \bottomrule
    \end{tabular}
\end{table}

\noindent
Example usage:
\begin{minted}{rust}
a = 3.7
b = a as i32     # b becomes 3
c = b as f32   # c becomes 3.0
\end{minted}


\subsubsection{Operator Precedence}
\begin{table}[H]
    \centering
    \begin{tabular}{|c|c|}
        \hline
        Operators & Associativity \\
        \hline
        Function Calls & \\
        Vector and String Subscripting & left-to-right \\
        Struct Member Access (.) & \\
        \hline
        Unary Operators (\mintinline{text}{+}, \mintinline{text}{-}, \mintinline{text}{~}) & \\
        Bitwise NOT & right-to-left \\
        Logical NOT & \\
        \hline
        Type Casting & left-to-right \\
        \hline
        \mintinline{text}{**} & right-to-left \\
        \hline
        \mintinline{text}{*, /, %, .*} & left-to-right \\
        \hline
        \mintinline{text}{+, -} & left-to-right \\
        \hline
        \mintinline{text}{<<, >>} & left-to-right \\
        \hline
        \mintinline{text}{&} & left-to-right \\
        \hline
        \mintinline{text}{^} & left-to-right \\
        \hline
        \mintinline{text}{|} & left-to-right \\
        \hline
        Comparison Operators & - \\
        \hline
        \mintinline{text}{&&} & left-to-right \\
        \hline
        Logical OR & left-to-right \\
        \hline
        Vector Concatenation & left-to-right \\
        \hline
        \mintinline{text}{=} & right-to-left \\
        \hline
        \mintinline{rust}{,} & left-to-right \\
        \hline
        \mintinline{rust}{return} & - \\
        \hline
    \end{tabular}
\end{table}

\subsection{Variable Declaration}
AFT supports type inference and explicit type annotations for variable declarations:
\begin{minted}[breaklines]{rust}
let a = 123;                      // i32 (inferred)
let b: u8 = 255;                  // explicit u8 type
let c: [c32] = [1, 2j, -3, -4j];  // complex array
let d: [c32] = [1, 2j, -3, -4j];  // same as above
let e = [1, 2j, -3, -4j];         // type inferred as [c64]
let f = 6 - 9j;                   // type inferred as c64

let complex_64_array = []; // ERROR: Type cannot be inferred

// Correct approach:
let complex_64_array: [c64] = [];

let x = 3.0;
x = 2; // ERROR: Assigning integer to float

let x = 3.0;
let x = 2; // WARNING: Variable x is shadowed

let a = 0+0j;  // Complex zero (c64)
let a = 0.;    // Float zero (f64)
let a = 0;     // Integer zero (i32)

let a = 1, b = 2.0, c = 4 + j ;
\end{minted}
Multiple variables can be declared in a single statement:
\begin{minted}{rust}
let (a,b,c) = (1,2.0,4+j) ;
let (a,b,a) = (1,2.0,4+j) ;         // ERROR : Attempt to shadow variable in same statement
\end{minted}
Parentheses are allowed on the left-hand side of variable declarations and assignment operators only if there are two or more comma-separated identifiers present.
Constants can be defined using the \mintinline{rust}{const} keyword. The type must be explicitly mentioned for constants.
\begin{minted}{rust}
// Definition of PI (a built-in constant)
const PI: f64 = 3.14159265358979323846264338327950288;
const TAU: f64 = 2. * PI;
const (PI_2, PI_4): (f64, f64) = (PI / 2., PI / 4.);
\end{minted}
The right-hand side of a constant definition must only have literals and other constant variables. Only \mintinline{rust}{const} variables can be global variables.

\subsection{Control Flows}

\subsubsection{Conditional Statements}
\begin{minted}{rust}
if condition {
    // do something
}

if condition {
    // do something
} else {

    // do something else
}

if condition1 {
    // do something
} else if condition2 {
    // otherwise do something
} else {
    // do something else
}
\end{minted}

\subsubsection{Loops}
\textbf{Repeat Loop}
\begin{minted}{rust}
repeat number_of_iterations {
    // do something
}
\end{minted}
\mintinline{text}{number_of_iterations} must be an unsigned integer. \\
\phantom{.} \\
\textbf{While Loop}
\begin{minted}{rust}
while condition {
    // do something
}
\end{minted}
\phantom{.} \\
\textbf{For Loop}
\begin{minted}{rust}
for element in collection {
    // do something
}
\end{minted}

\subsection{Functions}
Functions in AFT are defined as below:
\begin{minted}{rust}
    fn function_name(parameter_list) return_type_list {
        // do something
    }
\end{minted}
Calling functions is straightforward:
\begin{minted}{rust}
    function_name(param1, param2); // Example function call for function with 2 arguments
    let x = function_with_1_return(...);
    let (x, y) = function_with_2_return(...);
\end{minted}
Functions can return multiple values. Returned values must be enclosed in parentheses and separated by commas.
\begin{minted}{rust}
    fn sum_and_product(a: i32, b: i32) (i32, i32) {
        return (a + b, a * b);
    }
    let (x, y) = sum_and_product(3, 2);
\end{minted}
When assigning a function's returned values to variables, the number of variables must be equal to the number of returned values.
\begin{minted}{rust}
    (x, y) = sum_and_product(3, 2); // OK
    x = sum_and_product(3, 2); // ERROR: 
\end{minted}

\subsection{Structs}
AFT supports user-defined structs with members laid out in memory in the same order they are declared. These structs also support nesting.
\begin{minted}{rust}
// A simple struct with three 32-bit integers as members
struct Point3 {
    x: i32,
    y: i32,
    z: i32
}
\end{minted}
Structs can be initialized as follows:
\begin{minted}{rust}
// Create a Point3 instance with x = 1, y = 2 and z = 3.
let point_1 = Point3 { x: 1, y: 2, z: 3 };
\end{minted}
Struct members can be accessed using the \mintinline{rust}{.} operator.
\begin{minted}{rust}
let x_coordinate_of_point_1 = point_1.x;
\end{minted}

\subsection{Comments}
AFT supports both single-line and multi-line comments:

\begin{minted}{rust}
// This is a single line comment

/* 
This is a
multi-line
comment.
*/
\end{minted}
Single-line comments begin with \mintinline{rust}{//} and continue to the end of the line. Multi-line comments are enclosed between \mintinline{rust}{/*} and \mintinline{rust}{*/} and can span multiple lines.